\chapter{Revisão da literatura}
\label{cap:literatura}

Este capítulo situa o método adotado dentro da literatura mais ampla
sobre geração e uso de malhas de Voronoi, organizada em três grupos
temáticos: (i)~o núcleo dinâmico do próprio \mpasa{} e seu uso em
assimilação de dados; (ii)~geração e otimização de malhas de Voronoi
esféricas; e (iii)~interpolação sobre malhas de Voronoi/Delaunay fora do
domínio da geociência atmosférica, que ajudam a generalizar e comparar o
método baricêntrico adotado.

\section{O núcleo dinâmico MPAS-A e a malha SCVT}

A descrição de referência do núcleo dinâmico do \mpasa{}
\cite{skamarock2012voronoi} estabelece a malha $C$-\textit{grid} sobre
tesselação de Voronoi centroidal esférica como escolha estrutural do
modelo, motivada pela ausência de singularidades polares e pela
possibilidade de refinamento local suave -- exatamente as duas
propriedades que tornam a malha adequada tanto para a simulação quanto,
como demonstrado neste trabalho, para servir como origem/destino de uma
interpolação nativa sem grade intermediária. A apresentação de Duda
\cite{duda2014scalability} complementa essa descrição com considerações
de desempenho/escalabilidade da implementação de produção, relevantes
para justificar por que operações vetorizadas sobre a conectividade da
malha (como o cálculo de pesos baricêntricos célula a célula) são
tratáveis computacionalmente mesmo em malhas de centenas de milhares de
células.

O uso de interpolação baricêntrica na malha dual de Delaunay como método
de remapeamento entre estados do \mpasa{} já é prática estabelecida na
comunidade -- não uma proposta original deste trabalho, mas a aplicação
de um método já validado a um novo contexto (geração de condição
inicial/fronteira lateral, em vez de assimilação de dados). A referência
central para essa confirmação é Ha et al.
\cite{ha2017enkf}, que descreve o sistema de assimilação por
filtro de Kalman em conjunto (MPAS-DART) do NCAR: a Eq.~6 e a Fig.~1
daquele trabalho descrevem exatamente a interpolação baricêntrica no
triângulo dual de Delaunay (três centros de célula vizinhos) para campos
escalares de massa, e a Fig.~5 demonstra que a reconstrução vetorial do
vento nos centros de célula via Funções de Base Radial (RBF),
\emph{antes} de interpolar, produz resultado superior a tentar
interpolar a componente normal diretamente nas arestas -- é essa
constatação que justifica a ordem de operações adotada na
\cref{sec:vento_arestas} deste relatório. Como contraponto, os materiais
técnicos do sistema de assimilação MPAS-JEDI \cite{mpasjedi2021staticb}
revelam que mesmo abordagens mais recentes de assimilação variacional
ainda recorrem, para operações específicas sem inversa natural na malha
não estruturada (como a transformação $\psi,\chi \to u,v$ via
transformada espectral), a uma etapa intermediária em grade
latitude-longitude -- evidenciando que ``100\% nativo'' nem sempre é
possível para \emph{toda} operação de um sistema de modelagem/assimilação
completo, mas é possível, e é o padrão estabelecido, especificamente para
o remapeamento direto de campos entre duas malhas -- o problema atacado
neste trabalho.

\section{Geração e otimização de malhas de Voronoi esféricas}

A construção computacional de tesselações de Voronoi centroidais é, por
si, uma área de pesquisa ativa. Du e Gunzburger
\cite{dugunzburger2002cvt} formalizam a otimização de CVTs via
minimização de uma função energia relacionada ao algoritmo de Lloyd,
fundamento teórico do processo de geração da própria malha do \mpasa{}.
Yang, Gunzburger e Ju \cite{yang2017cvtsphere} propõem um método de
geração paralela mais rápido, baseado em um algoritmo LBFGS
pré-condicionado por Lloyd, especificamente para o caso esférico --
relevante para a geração de malhas regionais de alta resolução, embora
não usado diretamente na etapa de \emph{interpolação} deste trabalho (que
opera sobre malhas já geradas por outra ferramenta,
\texttt{MPAS-Limited-Area}). Engwirda \cite{engwirda2017jigsawgeo}
descreve o \textsc{jigsaw-geo}, um gerador de malha não estruturada
ortogonal localmente escalonada para modelagem de circulação geral,
alternativa de geração de malha usada por outros modelos da mesma
família dinâmica (por exemplo MPAS-Ocean); e Santos e Peixoto
\cite{santospeixoto2021refinement} exploram refinamento local de malha
esférica de Voronoi orientado por topografia em modelos de água rasa,
tema adjacente ao problema de suavização de terreno tratado na
\cref{cap:fase2}\footnote{Refinamento de \emph{malha} (mudar a resolução
espacial dos geradores conforme o relevo) e suavização do \emph{campo de
terreno} usado na definição da coordenada vertical (mantendo a malha
fixa, o problema efetivamente tratado neste trabalho) são technicamente
distintos, mas motivados pelo mesmo fenômeno físico -- erro de gradiente
de pressão horizontal sobre relevo íngreme.}.

\section{Interpolação sobre malhas de Voronoi/Delaunay em outros domínios}

Fora da modelagem atmosférica, a interpolação sobre malhas de Voronoi tem
tradição estabelecida em geometria computacional e em outras áreas da
física computacional, o que permite contextualizar o método baricêntrico
como caso particular de uma família mais ampla de interpolantes.
Hiyoshi e Sugihara \cite{hiyoshi2000continuity} apresentam um arcabouço
geral de interpolantes baseados no diagrama de Voronoi com continuidade
$C^1$ (generalizando os interpolantes clássicos de Sibson e de Laplace),
mais suaves que a interpolação baricêntrica linear por partes usada
neste trabalho -- um candidato natural de evolução futura caso a
descontinuidade de derivada nas arestas dos triângulos de Delaunay se
mostre relevante em aplicações futuras (\cref{cap:discussao}). Melodia e
Lenz \cite{melodia2019persistenthomology} propõem um critério de parada
baseado em homologia persistente para refinar iterativamente uma
interpolação de Voronoi, tema de otimização de malha não diretamente
aplicável ao caso deste trabalho (malhas fixas, pré-geradas), mas
ilustrativo da maturidade do campo. Em astrofísica e ciências da Terra,
Petkova et al.~\cite{petkova2017densitygridding} usam grades de Voronoi
para regularizar dados de hidrodinâmica lagrangiana (SPH), e Camps et
al.~\cite{camps2013radiativetransfer} e Bonduà et al.
\cite{bondua2017voronoi3d} aplicam malhas de Voronoi tridimensionais em
transferência radiativa astrofísica e migração de gás em formações
sedimentares, respectivamente -- ambos os casos usando a mesma
propriedade central explorada neste trabalho (adaptabilidade local de
resolução sem descontinuidade de malha), embora em três dimensões
espaciais em vez da superfície esférica bidimensional do \mpasa{}.
Abdelkader et al.~\cite{abdelkader2019vorocrust} tratam do problema
correlato, mas distinto, de geração de malha de Voronoi conformada a uma
superfície sem recorte (\emph{clipping}) dos poliedros de fronteira --
não diretamente aplicável aqui, já que a malha do \mpasa{} já vem
pré-gerada e recortada pela ferramenta \texttt{MPAS-Limited-Area}, mas
relevante como pano de fundo sobre os desafios numéricos de se trabalhar
com malhas de Voronoi em domínios limitados.

\section{Síntese}

A literatura revisada confirma dois pontos que sustentam as decisões de
projeto deste trabalho: primeiro, que a interpolação baricêntrica na
malha dual de Delaunay não é uma escolha \emph{ad hoc}, mas o método
padrão já validado pela comunidade MPAS (via MPAS-DART) para
remapeamento nativo de campos de massa entre estados dessa mesma malha;
segundo, que existem interpolantes de maior continuidade
(Hiyoshi--Sugihara) e técnicas mais sofisticadas de geração/refinamento
de malha que não foram necessários para o escopo deste trabalho -- que
lida com remapeamento entre duas malhas \emph{já geradas} -- mas que
constituem direções concretas de trabalho futuro caso a precisão do
método baricêntrico linear se mostre insuficiente em aplicações que
exijam maior suavidade de derivadas (por exemplo, esquemas de advecção
de ordem mais alta que dependam de gradientes horizontais contínuos do
campo interpolado).
